\begin{frame}
	\begin{center}
		\Huge Kampf
	\end{center}
\end{frame}

\begin{frame}
	\frametitle{Das Kampfkonzept}
		\begin{itemize}
		\item Klick auf den Lehrer $\rightarrow$ Kampf startet
		\item Schüler und Lehrer im Kampfhintergrund 
		\item Schüler-Angriff
		\item Lehrer verteidigt oder nicht
		\item Lehrer-Angriff
		\item Schüler verteidigt innerhalb von 4s ich änder eh nochmal
	\end{itemize}
\end{frame}

\begin{frame}
	\frametitle{Lösungsansätze}
	Erster Lösungsansatz: while-Schleife
	Hier kommt eine Aufzählung hin.
\end{frame}

\begin{frame}
	\frametitle{Funktion: getItems}
	\begin{figure}
		\centering
		\includegraphics[width=.8\textwidth]{bilder/GetItems.png}
		\caption{Ausschnitt der Funktion getItems (Bildschirmfoto des Projektes)}
		\label{fig:agi}
	\end{figure}
\end{frame}

\begin{frame}
	\frametitle{Lehrer-Funktion}
		\begin{figure}
		\centering
		\includegraphics[width=.6\textwidth]{bilder/Lehrer-Funktion.png}
		\caption{Ausschnitt der Funktion getItems (Bildschirmfoto des Projektes)}
		\label{fig:alf}
	\end{figure}
\end{frame}

\begin{frame}
	\frametitle{Angriff und Verteidigung}
	Das kommt dann noch.
\end{frame}

\begin{frame}
	Erst erklären warum Idee I nicht funktioniert.\\
	Dann Idee zwei erklären.\\
	\begin{itemize}
		\item Update Funktion noch einmal einführen
		\item warum sie uns hilft
		\item verschiedene updates erklären
		\item noch draw Funktion ansprechen
	\end{itemize}
\end{frame}

\begin{frame}
	\frametitle{Methode 2}
	\begin{figure}
		\centering
		\includegraphics[width=.9\textwidth]{bilder/InfoKampfMethode2.png}
		\caption{Warum Methode 1 nicht funktioniert}
		\label{fig:erklaerungMethode2}
	\end{figure}
\end{frame}

\begin{frame}
	\frametitle{Die update() Funktion}
	\begin{figure}
		\centering
		\includegraphics[width=.9\textwidth]{bilder/InfoUpdates.png}
		\caption{Wie update() funktioniert}
		\label{fig:erklaerungUpdates}
	\end{figure}
\end{frame}

\begin{frame}
	\begin{algorithm}[H]
		\begin{algorithmic}[1]
			\STATE Hintergrundbild
			\STATE Variablen festsetzen
			\STATE Lehrer $\gets$ new LehrerObject(FACHSCHAFT)
			\STATE button zum Kampfabbruch
			\STATE music $\gets$ Kampfmusik
		\end{algorithmic}
		\caption{KampfScene(FACHSCHAFT) pseudocode}
		\label{alg:SceneKonstruction}
	\end{algorithm}
\end{frame}

\begin{frame}
	\begin{algorithm}[H]
		\begin{algorithmic}[1]
			\IF {LehrerLeben $\leq$ 0}
				\STATE Gewonnen
				\STATE Dekrementiere UebrigeLehrer
				\STATE Beende Kampf
			\ENDIF
			\IF {PlayerLeben $\leq$ 0 $\vee$ |PlayerItems| = 0}
				\STATE Verloren
				\STATE Beende Kampf
			\ENDIF
		\end{algorithmic}
		\caption{KampfScene.update() pseudocode}
		\label{alg:SceneUpdate}
	\end{algorithm}
\end{frame}

\begin{frame}
	\begin{algorithm}[H]
		\begin{algorithmic}[1]
			\IF {timer $\neq$ 0}
			\STATE draw timer 
			\ENDIF
		\end{algorithmic}
		\caption{KampfScene.draw() pseudocode}
		\label{alg:SceneDraw}
	\end{algorithm}
\end{frame}

\begin{frame}
	\begin{algorithm}[H]
		\begin{algorithmic}[1]
			\IF {Turn = LehrerAngriff}
				\STATE ZwischenSchaden $\gets$ LehrerAngriff()
				\IF {ItemBenutzt} 
					\STATE DisplayItem()
				\ENDIF
				\STATE timer $\gets$ 4 Sekunden
				\STATE Turn $\gets$ PlayerVerteidigung
			\ELSIF {Turn = LehrerVerteidigung}
				\STATE LehrerLeben -= LehrerVerteidigung(Zwischenschaden)
				\STATE Turn $\gets$ PlayerAngriff
			\ENDIF
		\end{algorithmic}
		\caption{LehrerObject.update() pseudocode}
		\label{alg:LehrerUpdate}
	\end{algorithm}
\end{frame}

\begin{frame}
	\begin{algorithm}[H]
		\begin{algorithmic}[1]
			\IF {imKampf}
			\IF {Turn = PlayerAngriff}
			\STATE ZwischenSchaden $\gets$ PlayerAngriff(clickedItem)
			\STATE Turn $\gets$ LehrerVerteidigung
			\ELSIF {Turn = PlayerVerteidigung}
			\STATE PlayerLeben -= PlayerVerteidigung(Zwischenschaden, clickedItem)
			\STATE Turn $\gets$ LehrerAngriff
			\ENDIF
			\ENDIF
		\end{algorithmic}
		\caption{PlayerObject.update() pseudocode}
		\label{alg:PlayerUpdate}
	\end{algorithm}
\end{frame}

\begin{frame}
	\begin{algorithm}[H]
		\begin{algorithmic}[1]
			\IF {turn = Angriff}
			\STATE draw \dq Angriff"
			\ELSIF {turn = Verteidigung}
			\STATE draw "Verteidigung"
			\ENDIF
			\STATE draw Leben
		\end{algorithmic}
		\caption{draw() pseudocode}
		\label{alg:Draw}
	\end{algorithm}
\end{frame}
